\section{Discussion}
\label{sec:discussion}

\paragraph{Why is control nearly free?}
The properties are linear in composition, so guidance only re-weights first
moments at each position; it does not move the backbone-conditioned
high-density region. Consistent with this, the property-quality coupling we
measure \emph{between designs} is positive (GRAVY vs scTM $+0.61$, aggregation
vs scTM $+0.56$) but the range we move ($\Delta$GRAVY $\approx 0.4$) is small
relative to between-design dispersion, so no second-order collapse is
triggered.

\paragraph{When does it stop being free?}
Two observed boundaries: (i) pushing a property that is strongly coupled to
quality (GRAVY at $\lambda{=}4$ already lowers scTM to 0.9257, and
$\lambda{=}8$ yields unfolded sequences); (ii) pushing a property whose weight
vector is correlated with another (GRAVY $\to +20$ charge). Controllability is
therefore not a property of a model but of the (objective $\times$ weight
vector) pair.

\paragraph{Implications.}
For weakly coupled properties (surface hydrophobicity, net charge) retraining
is unnecessary: the closed-form tilt is larger in effect, reversible, and
cheaper. For strongly coupled objectives (stability, expression) a prior shift
or RL is still required. In both cases acceptance must report the properties
that were \emph{not} requested.
